%---------------------------------------------------------------
\chapter{\kesimpulan}
\label{bab:6}
%---------------------------------------------------------------
Pada bab ini, Penulis akan memaparkan kesimpulan penelitian dan saran untuk penelitian berikutnya.


%---------------------------------------------------------------
\section{Conclusion}
\label{sec:kesimpulan}
%---------------------------------------------------------------
This paper presents the digital visualization tool Latviz, which was created to help learners more easily understand the entire process behind the ML-KEM and ML-DSA cryptographic standards and the lattice basis reduction algorithms LLL and BKZ. As post-quantum cryptography becomes increasingly important, there is a growing need for educational resources that allow learners to understand not only the purpose of these algorithms but how they work. Existing resources often focus on high-level overviews or require substantial prior knowledge of cryptography. To address this gap, Latviz was developed to provide complete step-by-step visualizations for ML-KEM and ML-DSA, accompanied by supplementary explanations of the functions, variables, and inner processes involved in their execution, and algorithm tracing complete with explanations for LLL and BKZ.

The resulting tool enables users to explore all four algorithms interactively, inspect intermediate values and operations, and follow the transformation of inputs into outputs throughout the execution process. By combining detailed visualizations with explanatory content, Latviz aims to make complex post-quantum cryptographic concepts more accessible to learners with varying levels of prior cryptographical knowledge.

To evaluate the effectiveness of the tool, a survey involving 22 participants sampled through convenience sampling with differing levels of familiarity with cryptography and post-quantum cryptography was conducted. The results indicate that Latviz was generally perceived as an effective learning tool for both algorithms, with 77.3\%, 81.8\%, 72.7\%, and 77.3\% of respondents considering it effective for ML-KEM, ML-DSA, LLL, and BKZ respectively. Furthermore, 81.8\% of respondents reported that the tool improved their understanding of ML-KEM and ML-DSA, and 72.7\% reported the same for LLL and BKZ. Relating to ML-KEM and ML-DSA, participants responded particularly positively to the step-by-step visualizations and their ability to illustrate the relationships between algorithmic processes. However, the results also suggest that users with limited prior knowledge of cryptography may have trouble with following the provided explanations and may benefit from additional context for unfamiliar concepts. For LLL and BKZ, participants found the calculation details panel and step-by-step visualizations to be the most helpful, while the sections relegated to explaining the roles of variables were found to be least effective. 

Overall, the evaluation results indicate that Latviz serves as an effective tool for improving users' understanding of ML-KEM, ML-DSA, LLL, and BKZ. These findings demonstrate the potential of interactive visualization as an educational approach for helping learners engage with and understand complex post-quantum cryptographic algorithms.
%---------------------------------------------------------------
\section{Future Works}
\label{sec:saran}
%---------------------------------------------------------------
While the evaluation results indicate that Latviz is effective in supporting users’ understanding of ML-KEM, ML-DSA, LLL, and BKZ, there are still several opportunities for improvement. The survey responses suggest that some users, particularly those with limited prior knowledge of cryptography, found certain explanations difficult to follow. Future development could therefore focus on improving the accessibility of the content through the addition of or references to introductory materials, more contextual explanations, and alternative methods of presenting complex concepts.
Future work may also involve expanding the scope of the tool by incorporating additional post-quantum cryptography material such as the NTRU cryptosystem. Additionally, evaluations involving larger and more diverse participant groups as well as pre-test and post-test assessments could provide a more precise and comprehensive measurement of the tool’s effectiveness in improving users’ understanding of post-quantum cryptographic concepts. Aside from technical improvements, increasing support for different languages can also improve the tool for audiences more comfortable with different languages.

